% Единственный источник содержания темы. UTF-8.
% Метки label и ключи в [...] сохраняют опубликованные ссылки.

\section{Теория}

\subsection{Определение}

На человеческом языке события называются \emph{независимыми}, если вероятность наступления первого события не зависит от вероятности наступления второго, и наоборот вероятность второго не зависит от первого. Но вам надо знать строгие определения ниже.

\begin{definition}
\emph{Первое определение.} События $A$ и $B$ называются \textbf{независимыми}, если

\[\mathbb{P}(A\cap B) = \mathbb{P}(A)\cdot \mathbb{P}(B)\]
\end{definition}

\begin{definition}
\emph{Второе определение.} При $\mathbb{P}(B)>0$ события $A$ и $B$ называются \textbf{независимыми}, если $\mathbb{P}(A\vert B) = \mathbb{P}(A)$.
\end{definition}

\begin{task}[independent-events/theory/1]\label{task-theory-1}
\textbf{Задача 1.} Докажите, что определения эквивалентны при $\mathbb{P}(B)>0$.
\end{task}

\begin{solution}[independent-events/theory/1]\label{solution-independent-events-theory-1}
Пусть $\mathbb P(B)>0$. Если выполнено первое определение, то

\[\mathbb P(A\mid B)=\frac{\mathbb P(A\cap B)}{\mathbb P(B)}=\frac{\mathbb P(A)\mathbb P(B)}{\mathbb P(B)}=\mathbb P(A).\]

Получили второе определение. Наоборот, из равенства $\mathbb P(A\mid B)=\mathbb P(A)$ и формулы условной вероятности после умножения на $\mathbb P(B)$ следует $\mathbb P(A\cap B)=\mathbb P(A)\mathbb P(B)$. Значит, при указанном условии определения эквивалентны.
\end{solution}

\begin{task}[independent-events/theory/2]\label{task-theory-2}
\textbf{Задача 2.} Докажите, что из второго определения также почти следует, что $\mathbb{P}(B) = \mathbb{P}(B\vert A)$. Приведите пример, когда это неверно.
\end{task}

\begin{solution}[independent-events/theory/2]\label{solution-independent-events-theory-2}
Если дополнительно $\mathbb P(A)>0$, то из независимости получаем

\[\mathbb P(B\mid A)=\frac{\mathbb P(A\cap B)}{\mathbb P(A)}=\mathbb P(B).\]

Оговорка нужна из-за возможности $\mathbb P(A)=0$. Например, при броске кубика возьмём $A=\varnothing$, $B=\{1\}$. Тогда $\mathbb P(B)=\frac16>0$, а $\mathbb P(A\mid B)=0=\mathbb P(A)$: второе определение выполнено. Однако $\mathbb P(B\mid A)$ не определена, поскольку пришлось бы делить на ноль.

Итак, при нулевой вероятности $A$ обратное условное равенство нельзя записывать; это не нарушение симметрии независимости, а ограничение определения условной вероятности.
\end{solution}

\begin{task}[independent-events/theory/3]\label{task-theory-3}
\textbf{Задача 3.} Существует ли событие $A$ независимое с любым другим событием?
\end{task}

\begin{solution}[independent-events/theory/3]\label{solution-independent-events-theory-3}
Да. Например, $A=\Omega$ независимо от любого события $B$:

\[\mathbb P(\Omega\cap B)=\mathbb P(B)=1\cdot\mathbb P(B)=\mathbb P(\Omega)\mathbb P(B).\]

Ещё один пример — $A=\varnothing$: обе стороны равенства независимости равны нулю.
\end{solution}

\begin{theorem}[independent-events/theorem/1]\label{theorem-1}
Если $A$ и $B$ независимы, то каждая из следующих пар событий тоже независима: $A, \overline{B};\ \overline{A},B;\ \overline{A}, \overline{B}$
\end{theorem}

\begin{proof}[independent-events/theorem/1]\label{solution-independent-events-theorem-1}
По независимости $\mathbb P(A\cap B)=\mathbb P(A)\mathbb P(B)$. Событие $A$ разбивается на несовместные части $A\cap B$ и $A\cap\overline B$. Поэтому

\[\mathbb P(A\cap\overline B)=\mathbb P(A)-\mathbb P(A\cap B)=\mathbb P(A)(1-\mathbb P(B))=\mathbb P(A)\mathbb P(\overline B).\]

Значит, $A$ и $\overline B$ независимы. Аналогично,

\[\mathbb P(\overline A\cap B)=\mathbb P(B)-\mathbb P(A\cap B)=(1-\mathbb P(A))\mathbb P(B).\]

Наконец, $\overline A\cap\overline B$ противоположно $A\cup B$. По формуле сложения

\[\mathbb P(\overline A\cap\overline B)=1-\mathbb P(A)-\mathbb P(B)+\mathbb P(A\cap B).\]

Подставляя независимость, получаем $(1-\mathbb P(A))(1-\mathbb P(B))=\mathbb P(\overline A)\mathbb P(\overline B)$. Все три пары независимы. В доказательстве нет деления на вероятности, поэтому оно работает и при вероятностях 0 или 1.
\end{proof}

\section{Задачи}

\textbf{!!!} ОПИРАТЬСЯ НА ИНТУИЦИЮ НЕЛЬЗЯ \textbf{!!!}\\

\begin{task}[independent-events/problems/1]\label{task-problems-1}
\textbf{4.} Рассмотрим двукратное бросание игральной кости и два события: событие А «в первый раз выпало более трёх очков» и событие В «во второй раз выпало менее трёх очков». Будут ли события А и В независимыми?
\end{task}

\begin{solution}[independent-events/problems/1]\label{solution-independent-events-problems-1}
Всего 36 равновозможных пар результатов. Для события $A$ первый результат принадлежит $\{4,5,6\}$, поэтому $\mathbb P(A)=\frac36=\frac12$. Для события $B$ второй результат принадлежит $\{1,2\}$, поэтому $\mathbb P(B)=\frac26=\frac13$.

Пересечению благоприятствуют $3\cdot2=6$ пар, откуда

\[\mathbb P(A\cap B)=\frac6{36}=\frac16=\frac12\cdot\frac13=\mathbb P(A)\mathbb P(B).\]

\textbf{Ответ:} события независимы.
\end{solution}

\begin{task}[independent-events/problems/2]\label{task-problems-2}
\textbf{5.} Кубик вновь бросают два раза. Событие А "сумма результатов бросков кратна 3" и событие В "сумма результатов бросков чётна". Зависимы ли эти события?
\end{task}

\begin{solution}[independent-events/problems/2]\label{solution-independent-events-problems-2}
Проверим определение независимости по 36 равновозможным парам результатов.

Сумма кратна 3, если она равна 3, 6, 9 или 12. Число соответствующих пар равно $2+5+4+1=12$, поэтому $\mathbb P(A)=\frac13$.

Сумма чётна, когда оба результата одной чётности: $3\cdot3+3\cdot3=18$ пар. Поэтому $\mathbb P(B)=\frac12$.

В пересечении сумма должна быть и чётной, и кратной 3, то есть равной 6 или 12. Таких пар $5+1=6$:

\[\mathbb P(A\cap B)=\frac6{36}=\frac16=\mathbb P(A)\mathbb P(B).\]

\textbf{Ответ:} события независимы.
\end{solution}

\begin{task}[independent-events/problems/3]\label{task-problems-3}
\textbf{6.} Всем пациентам с подозрением на гепатит делают анализ крови. Если анализ выявляет гепатит, то результат анализа называется положительным. У больных гепатитом пациентов анализ даёт положительный результат с вероятностью $0,9$. Если пациент не болен гепатитом, то анализ может дать ложный положительный результат с вероятностью $0,01$. Известно, что $5\%$ пациентов, поступающих с подозрением на гепатит, действительно больны гепатитом. Найдите вероятность того, что результат анализа у пациента, поступившего в клинику с подозрением на гепатит, будет положительным.
\end{task}

\begin{solution}[independent-events/problems/3]\label{solution-independent-events-problems-3}
Обозначим через $H$ наличие гепатита, через $T$ — положительный результат анализа. События $H$ и $\overline H$ разбивают всех пациентов на две несовместные группы.

Вероятность положительного анализа у больного пациента:

\[\mathbb P(H\cap T)=0,05\cdot0,9=0,045.\]

Вероятность ложного положительного результата у здорового:

\[\mathbb P(\overline H\cap T)=0,95\cdot0,01=0,0095.\]

Складываем вероятности двух несовместных способов получить положительный результат:

\[\mathbb P(T)=0,045+0,0095=0,0545.\]
\end{solution}

\begin{task}[independent-events/problems/4]\label{task-problems-4}
\textbf{7.[ЕГЭ]} Вероятность того, что батарейка бракованная, равна $0,06$. Покупатель в магазине выбирает случайную упаковку, в которой две таких батарейки. Найдите вероятность того, что обе батарейки окажутся исправными. Считайте исправность батареек независимой.
\end{task}

\begin{solution}[independent-events/problems/4]\label{solution-independent-events-problems-4}
Вероятность исправности одной батарейки равна $1-0,06=0,94$. По условию исправности двух батареек независимы. Поэтому

\[\mathbb P(\text{обе исправны})=0,94\cdot0,94=0,8836.\]
\end{solution}

\begin{task}[independent-events/problems/5]\label{task-problems-5}
\textbf{8.[ЕГЭ]} В магазине три продавца. Каждый из них занят с клиентом с вероятностью 0,3. Найдите вероятность того, что в случайный момент времени все три продавца свободны одновременно (считайте, что клиенты заходят независимо друг от друга).
\end{task}

\begin{solution}[independent-events/problems/5]\label{solution-independent-events-problems-5}
Каждый продавец свободен с вероятностью $1-0,3=0,7$. Используем предусмотренную в задаче независимую модель занятости продавцов. Вероятность одновременной свободы всех троих равна

\[0,7\cdot0,7\cdot0,7=0,343.\]
\end{solution}

\begin{task}[independent-events/problems/6]\label{task-problems-6}
\textbf{9.[ЕГЭ]} Помещение освещается фонарём с двумя лампами. Вероятность перегорания лампы в течение года равна 0,3. Найдите вероятность того, что в течение года хотя бы одна лампа не перегорит. Считайте перегорание ламп независимым.
\end{task}

\begin{solution}[independent-events/problems/6]\label{solution-independent-events-problems-6}
Противоположное событие — обе лампы перегорят за год. По условию перегорания независимы, поэтому вероятность этого события равна $0,3^2=0,09$.

\textbf{Ответ:} вероятность, что хотя бы одна лампа не перегорит, равна $1-0,09=0,91$.
\end{solution}

\section{Домашнее задание}

Вторая часть стр. 44-47.

\begin{task}[independent-events/homework/1]\label{task-homework-1}
\textbf{1.} Две фабрики выпускают одинаковые стекла для автомобильных фар. Первая фабрика выпускает 45\% этих стекол, вторая  —  55\%. Первая фабрика выпускает 3\% бракованных стекол, а вторая  —  1\%. Найдите вероятность того, что случайно купленное в магазине стекло окажется бракованным.
\end{task}

\begin{solution}[independent-events/homework/1]\label{solution-independent-events-homework-1}
Стекло может оказаться бракованным двумя несовместными способами: его изготовила первая фабрика или вторая. По формуле умножения вероятности этих случаев равны $0,45\cdot0,03$ и $0,55\cdot0,01$.

Складываем:

\[\mathbb P(\text{брак})=0,45\cdot0,03+0,55\cdot0,01=0,0135+0,0055=0,019.\]
\end{solution}

\begin{task}[independent-events/homework/2]\label{task-homework-2}
\textbf{2.} В магазине стоят два платёжных автомата. Каждый из них может быть неисправен с вероятностью 0,05 независимо от другого автомата. Найдите вероятность того, что хотя бы один автомат исправен.
\end{task}

\begin{solution}[independent-events/homework/2]\label{solution-independent-events-homework-2}
Противоположное событие — неисправны оба автомата. Неисправности независимы, поэтому его вероятность равна $0,05^2=0,0025$.

\textbf{Ответ:} $1-0,0025=0,9975$.
\end{solution}

\begin{task}[independent-events/homework/3]\label{task-homework-3}
\textbf{3.} Вероятность того, что телефон не сломается в течение первого года использования равна $0,97$, вероятность того, что не сломается в течение первых двух лет использования равна $0,87$. Найдите вероятность того, что телефон сломается не раньше первого года и не позже второго. (считаем, что вероятность поломки телефона ровно на первый или на второй или на третий или ... год с точностью до наносекунды равна 0)
\end{task}

\begin{solution}[independent-events/homework/3]\label{solution-independent-events-homework-3}
Пусть $A$ — телефон прослужит больше одного года, $B$ — больше двух лет. Тогда $B\subseteq A$. Искомое событие — $A\setminus B$: телефон пережил первый год, но не пережил второй.

Разбиение $A=(A\setminus B)\cup B$ состоит из несовместных событий, поэтому

\[\mathbb P(A\setminus B)=\mathbb P(A)-\mathbb P(B)=0,97-0,87=0,10.\]

Указанная в условии нулевая вероятность поломки точно на границе года позволяет не различать строгие и нестрогие границы периода.
\end{solution}

\begin{task}[independent-events/homework/4]\label{task-homework-4}
\textbf{4.} Найдите ошибку в решении следующей задачи:\\
 \textbf{Условие:} Вероятность того, что телефон не сломается в течение первого года использования равна $0,97$, вероятность того, что не сломается в течение первых двух лет использования равна $0,87$. Найдите вероятность того, что телефон сломается не раньше первого года и не позже второго.\\
 

\begin{given}[independent-events/homework/4]\label{given-solution-independent-events-homework-4}
\textbf{Проверяемое решение:} Вероятность того, что телефон сломается быстрее чем за 2 года равна $1 - 0,87 = 0,13$. Тогда вероятность того что телефон проживёт больше одного года и меньше двух равна $0,97\cdot 0,13 = 0,1261$.
\end{given}
\end{task}

\begin{solution}[independent-events/homework/4]\label{solution-independent-events-homework-4}
В предложенном решении перемножены вероятности событий «переживёт первый год» и «сломается до конца второго года» без основания считать эти события независимыми. Формулу произведения в таком виде применять нельзя.

Обозначим через $A$ выживание после первого года, через $B$ — после второго. Поскольку $B\subseteq A$, искомая вероятность равна

\[\mathbb P(A\setminus B)=0,97-0,87=0,10.\]

Можно проверить ошибку и через условную вероятность: вероятность поломки до конца второго года при условии выживания после первого равна $\frac{0,10}{0,97}=\frac{10}{97}$, а не $0,13$. Правильное произведение даёт $0,97\cdot\frac{10}{97}=0,10$.
\end{solution}

